\section{What a sparse morning does not determine}
\label{sec:infflat}

Chapter~9 closed with a moving prior: a surface transported to the
new spot, waiting to be compared with fresh evidence.  Before
building the update that weighs the two, it pays to ask a static
question --- one that needs no dynamics at all.  Take one morning's
quotes, alone.  How much of the smile do they actually determine?

The book's running node answers concretely.  The frozen SPY December
2026 chain carries \MacPriorFlatNFull\ prepared quotes
(\MacPriorFlatHeroDays\ days to expiry).  Keep only the
\MacPriorFlatNThin\ quotes inside the at-the-money band
$\lvert k\rvert\le0.10$ and set the rest aside.  On an illiquid
name, or in the first minutes of a session, a band like this is all
the market offers, and the desk still needs a whole smile --- the
wings included --- to price, hedge, and feed every integral of
Chapter~5.

What can the band say about the wing?  Here is a direct experiment,
run with a deliberately flexible curve family --- a \emph{spline},
which is just a smooth curve pieced together from cubic polynomials
and controlled by its values at a handful of chosen strikes (the
\emph{knots}; the exact knots and the small smoothness penalty are
in Appendix~\ref{app:infprotocol}) --- so that no parametric habit
answers in the data's place.  Fit the \MacPriorFlatNThin\ band quotes, but
\emph{impose} the wing: pin the fitted volatility at $k=-0.30$ to a
chosen value, and refit.  Sweep the pinned value from
\MacPriorFlatPinLoPct\% to \MacPriorFlatPinHiPct\% --- a span of
\MacPriorFlatPinSpanPts\ volatility points --- and record, for each
imposed wing, how well the curve fits the quotes it was actually
given.  Each quote's misfit is measured against that quote's own
stated noise, the bid--ask half-spread; ``fits the band'' always
means ``fits within the noise the market itself states.''

\begin{figure}[htbp]
\centering
\paperfig{\textwidth}{fig_flt_flat.pdf}
\caption{What a sparse morning does not determine (frozen SPY
December 2026 node, \MacPriorFlatNThin\ quotes kept in the shaded
band).  (a)~Five fits of the same band, each with the wing volatility
at $k=-0.30$ (dotted vertical) pinned to a different imposed value;
through the band they are one curve.  (b)~The root-mean-square quote
error against the imposed wing value: on the thinned morning
(orange) it varies by \MacPriorFlatValleyBp\ vol bp across the whole
twelve-point sweep --- flat; on the full \MacPriorFlatNFull-quote
chain (ink) the same sweep meets a sharp V.  (c)~The quote-support
profile $\mathcal{Q}(k)$ of \cref{eq:infsupport} for both quote sets:
the thinned morning's support collapses outside the band, exactly
where panel~(b)'s valley is flat.}
\label{fig:infflat}
\end{figure}

\Cref{fig:infflat} shows the result.  In panel~(a), the orange dots are the morning's quotes and the five
curves are five fits, each pinned to a different wing value ---
$26\%$ up to $34\%$ at $k=-0.30$.  Through the shaded band the five
curves lie on top of one another; to the left they fan apart by
eight volatility points.  Panel~(b) makes the visual precise.  The
orange line is the fitting error of the band quotes, in vol bp, as
the imposed wing sweeps across its twelve-point range: it moves by
\MacPriorFlatValleyBp\ vol bp in total --- the band cannot tell the
imposed wings apart.  The ink line repeats the identical sweep on
the \emph{full} chain, where real quotes sit at $k=-0.30$: now the
error climbs steeply away from one preferred value, and an imposed
wing only half a volatility point off already costs five bp of
fitting error (\MacPriorFlatVHalfPts\ points is the half-width at
that cost).  Same family, same sweep, same objective --- the only
difference is which quotes exist.

This deserves a name.  Think of the fit's objective --- the sum of
squared, noise-scaled quote errors --- as a landscape over all
candidate curves.  (Statisticians call the corresponding probability
of the data the \emph{likelihood}; for Gaussian noise the objective
is its negative logarithm, so ``the likelihood is flat'' and ``the
objective is flat'' say the same thing, and the chapter uses both.)  A \emph{flat direction} is a way of changing the
curve that does not change the objective at all: the data assigns
the same score to every curve along it.  Panel~(b)'s orange line
exhibits one flat direction explicitly (move the wing, hold the
band), and there are many more --- every reshaping of the smile
outside the band, and finer wiggles between quotes inside it.
Statisticians say the wing is \emph{unidentified}: no amount of
reweighting or clever optimization extracts from these quotes a
number the quotes do not contain.

Two remarks before anything is built on this.

First, agreement between models is not evidence.  Refit the same
thinned band with the reference implementation's three parametric
families at several resolutions and smoothing strengths ---
eight fits in all --- and their wing readings at $k=-0.30$ land
within \MacPriorFlatEnsembleSpreadPts\ volatility points of one
another.  Eight rules agreeing looks like knowledge; panel~(b) says
the quotes would have accepted any answer across a twelve-point
range.  The agreement reflects a shared habit of smooth
extrapolation, a convention baked into every one of those families
--- not information.  A desk that reads consensus wings as measured
wings is mistaking its tools for its data.

Second, flatness is measurable in advance, cheaply.  The profile
experiment is exact but costs a refit per candidate value.  A proxy
does almost as well: count how many quotes sit \emph{near} a given
strike.  Define the \emph{quote support} at log-moneyness $k$ as
the kernel-weighted count --- a \emph{kernel} being nothing more
than a smooth bump, here a Gaussian, that decides how much a quote
at some distance still counts ---
\begin{equation}
\mathcal{Q}(k)
=\sum_{i}\exp\!\left(-\frac{(k_i-k)^2}{2\cdot0.06^2}\right),
\label{eq:infsupport}
\end{equation}
where the sum runs over the morning's quote locations $k_i$ (when
quotes carry unequal fit weights, each term is multiplied by its
weight, normalized to average one).  A quote sitting exactly at $k$
contributes one; a quote a few bandwidths of $0.06$ away contributes
essentially nothing.  Panel~(c) plots $\mathcal{Q}$ for both quote
sets on a logarithmic axis: the full chain holds double-digit
support everywhere the book has ever read that smile, while the
thinned morning's support falls off a cliff at the band edge ---
$\mathcal{Q}(0)=\MacPriorFlatSuppAtm$ at the money,
$\mathcal{Q}(-0.30)=\MacPriorFlatSuppWing$ in the wing.  Where
support is large the profile is a V; where it vanishes the profile
is flat.  Support is a proxy, not a theorem --- \cref{sec:infprice}
returns to exactly where it fails --- but it prices the geometry of
panel~(b) without a single refit.

We now know precisely what is missing: along the flat directions
the data has no opinion, and yet a published smile must take a
value there.  \emph{Something} always chooses --- a family's tail
habit, a smoothness penalty, a default.  The thesis of this book
calls such a choice a convention, to be stated as one, and the
natural candidate is obvious: yesterday's surface, which was
fitted to a full chain and has merely aged a day.  The rest of the
chapter builds the two instruments that let yesterday speak
\emph{only} where today is silent or noisy.  \Cref{sec:infgate}
through \cref{sec:infcoords} work in parameter space: where may a
prior pull at all?  \Cref{sec:inffilter} through \cref{sec:infaudit}
work in time: when today's data is present but noisy, how hard may
yesterday's prediction resist it?
